% =============================================================================
% Chapter 1 - Introduction
% =============================================================================
% Included from main.tex via % =============================================================================
% Chapter 1 - Introduction
% =============================================================================
% Included from main.tex via % =============================================================================
% Chapter 1 - Introduction
% =============================================================================
% Included from main.tex via % =============================================================================
% Chapter 1 - Introduction
% =============================================================================
% Included from main.tex via \include{contents/introduction}
%
% Cross-references used below (Structure of the Thesis only):
%   ch:literature            -> Chapter 3 (Literature Review)
%   sec:synthesis-and-gap    -> Section 3.10  (Synthesis and Gap)
%
% Labels defined here: ch:introduction, sec:motivation, sec:objectives,
%   sec:intro-contributions, sec:structure
%
% Bibliography keys: yan2014 (Yan et al., CASME II), xia2020a (Xia et al., STRCN)
%
% The aim, research question and objectives in Section 1.2 follow the third
% paragraph of the task description (contents/task_description.tex).
% The figures in Section 1.3 are taken from Chapter 5 (sec:transformer-effect,
% sec:evm-effect-result, sec:simam-effect-result, sec:cnn-effect-result,
% sec:headline-result) and must be kept consistent with it.
% =============================================================================

\chapter{Introduction}
\label{ch:introduction}

Micro-expressions are brief, involuntary facial movements that reveal an emotion which a person is attempting to conceal~\cite{yan2014}. They last no longer than approximately half a second~\cite{yan2014}, and their intensity is so low that an observer watching in real time will usually fail to notice them. A micro-expression may consist, for example, of a slight tightening of the brow or a movement at one corner of the mouth that is inconsistent with the expression the person is deliberately displaying.

This combination of informative content and near-invisibility has motivated sustained research into their automatic recognition. Yan et al.~\cite{yan2014} motivate the CASME II dataset on the grounds that a robust automatic recogniser ``would have broad applications in national safety, police interrogation, and clinical diagnosis''; the capability is likewise reported as promising for lie detection, business negotiation and psychoanalysis~\cite{xia2020a}. These are the motivating claims of the field rather than records of deployed systems, but they account for the continued research interest in the problem.

Automatic recognition systems for this task are typically assembled as pipelines of techniques that were introduced to the field separately. This thesis decomposes one such pipeline and measures the contribution of each of its stages individually.


\section{Research Motivation}
\label{sec:motivation}

Automatic recognition of micro-expressions from video cannot be achieved by applying a classifier directly to face images. Three properties of the phenomenon work against it. The facial movement is short, so only a small number of video frames carry the information that distinguishes one emotion from another. It is faint, so the information those frames carry lies close to the noise level of ordinary video. And it cannot be produced on instruction, so the recordings that contain it must be elicited under controlled laboratory conditions, which limits the resulting datasets to a small size.

Systems built for the task have responded to these constraints in a consistent way. The small facial motion is amplified, the sequence of pixel intensities is re-expressed as a description of movement, and the clip is resampled to a fixed duration before any learned component receives it; a spatial feature extractor then reduces each frame, and a temporal model processes the resulting sequence. Each stage was introduced with its own justification and its own reported gain, and each tends to be adopted together with the others.

This arrangement presents a practical difficulty for anyone building such a system. Implementation effort, and the tuning of preprocessing parameters that have no principled setting, are finite, so the developer must decide which stages justify that investment. The published literature does not answer this question, because each component is almost always evaluated within a complete system that differs from its predecessor in more than one respect.


\section{Research Aim and Objectives}
\label{sec:objectives}

The aim of this thesis is to investigate whether improved temporal modelling and motion representations can enhance the recognition of micro-expressions in dynamic facial expression recognition. The central research question is accordingly: \textit{how can the modelling of temporal phases and motion cues improve recognition accuracy for subtle and short-lived facial expressions?} The intended outcome is a clearer understanding of the limitations of current approaches, and of how sensitivity to fine-grained facial dynamics can be improved while robustness to noise is maintained.

The work builds on established supervised models and evaluates practical strategies for temporal and motion modelling. It does so by treating a standard micro-expression recognition pipeline as the object of study and measuring what each of its stages separately contributes under a single fixed protocol. Every configuration receives the same hybrid motion representation, in which each clip is encoded as horizontal optical flow, vertical optical flow and optical strain. Four components are then varied: Eulerian video magnification, which amplifies subtle facial motion before the motion representation is computed; a convolutional spatial stem, which learns spatial features from that representation; parameter-free attention (SimAM), which re-weights the learned features; and a temporal transformer, which models the sequence through temporal self-attention. The central research question is therefore addressed through three more specific questions:

\begin{enumerate}
	\item To what extent does temporal modelling with a transformer encoder improve recognition, compared with an otherwise identical pipeline without it?
	\item Does amplifying subtle facial motion through Eulerian video magnification, before the motion representation is computed, improve recognition?
	\item Do a learned convolutional spatial stem and parameter-free attention add to the recognition performance obtained from the hand-computed motion representation alone?
\end{enumerate}

To answer these questions, the thesis pursues the following objectives:

\begin{enumerate}
	\item \textbf{To implement a modular recognition pipeline} in which each of the four components can be enabled or disabled independently, with every disabled component replaced by a near-parameter-free fallback so that a comparison removes the component and nothing else.
	\item \textbf{To train and evaluate every architecturally valid configuration under one protocol.} The four components, switched on and off independently, give sixteen possible combinations, of which twelve are architecturally valid. Each of the twelve is trained from scratch, without pre-training or auxiliary data, which is the regime that a single small dataset permits. Each is evaluated under complete 25-fold leave-one-subject-out cross-validation on CASME II, using 156 clips in three classes of 99, 32 and 25 examples once the residual \texttt{Others} category is excluded. Because one fold is held out per subject, no subject's clips appear in both the training and the test data. Every factor not under variation is held identical across the twelve runs, and all twelve are scored by the same primary metric, pooled macro F1.
	\item \textbf{To estimate the marginal contribution of each component} from matched pairs of configurations. Any two configurations that differ in exactly one component isolate that component with all other factors fixed, so the difference between their scores estimates its contribution.
	\item \textbf{To examine the limitations of the resulting models}, by analysing per-class behaviour on the minority classes, comparing the results with published work, and stating the conditions under which the measured effects hold.
\end{enumerate}


\section{Research Contributions}
\label{sec:intro-contributions}

This thesis makes the following contributions.

\begin{enumerate}
	\item \textbf{A matched-pair factorial evaluation of a standard micro-expression recognition pipeline.} The reviewed literature typically evaluates components within complete systems, so that an improvement can rarely be attributed to the component credited with it. This thesis evaluates all twelve architecturally valid combinations of four components on CASME II under an identical leave-one-subject-out protocol, so that the contribution of each component, including learned ones, can be measured against an otherwise identical pipeline.

	\item \textbf{A measured account of the contribution of temporal and motion modelling.} Temporal modelling carries almost all of the measured effect: the temporal transformer improves pooled macro F1 by $+0.2173$ on average, and all six of its matched pairs are positive. The remaining components contribute far less. Eulerian video magnification has a small mean effect ($+0.0152$) whose direction varies across pairs; parameter-free attention has a mean effect ($+0.0034$) that cannot be distinguished from noise at this sample size; and the convolutional spatial stem is the only component with a negative mean effect ($-0.0310$).

	\item \textbf{Evidence that the full pipeline is not the best configuration.} A temporal transformer applied directly to the motion representation achieves a pooled macro F1 of $0.7122$, compared with $0.6659$ for the complete configuration in which all four components are enabled. An evaluation of the complete system alone would have reported the lower of these figures, with no means of discovering that a subset of its own architecture performs better.
\end{enumerate}


\section{Structure of the Thesis}
\label{sec:structure}

The thesis is organised in six chapters.

\begin{itemize}
	\item \textbf{Chapter 1: Introduction} motivates the study, states its aim, research questions and objectives, and summarises its contributions.

	\item \textbf{Chapter 2: Background} introduces the concepts on which the thesis builds: micro-expressions and the datasets used to study them, including CASME~II; data acquisition and preprocessing; Eulerian video magnification; optical flow and optical strain; neural network fundamentals and building blocks; training with limited and imbalanced data; and evaluation methodology.

	\item \textbf{Chapter 3: Literature Review} reviews the published evidence for each pipeline component and closes with the consolidated research gap and the declared divergences of the implemented system from the reviewed methods.

	\item \textbf{Chapter 4: Methodology} specifies the study as carried out: the experimental design, the dataset and its preparation, the twelve configurations, the architecture and its ablatable components, the training procedure and the evaluation protocol, together with the research gaps and limitations of each motion- and temporal-modelling component.

	\item \textbf{Chapter 5: Results} reports the results of the ablation, taking each component in turn over the matched pairs that isolate it, and then examines per-class behaviour, compares the results with published work and discusses them.

	\item \textbf{Chapter 6: Conclusion} answers the research questions and states the contributions, limitations and future work of the thesis.
\end{itemize}

%
% Cross-references used below (Structure of the Thesis only):
%   ch:literature            -> Chapter 3 (Literature Review)
%   sec:synthesis-and-gap    -> Section 3.10  (Synthesis and Gap)
%
% Labels defined here: ch:introduction, sec:motivation, sec:objectives,
%   sec:intro-contributions, sec:structure
%
% Bibliography keys: yan2014 (Yan et al., CASME II), xia2020a (Xia et al., STRCN)
%
% The aim, research question and objectives in Section 1.2 follow the third
% paragraph of the task description (contents/task_description.tex).
% The figures in Section 1.3 are taken from Chapter 5 (sec:transformer-effect,
% sec:evm-effect-result, sec:simam-effect-result, sec:cnn-effect-result,
% sec:headline-result) and must be kept consistent with it.
% =============================================================================

\chapter{Introduction}
\label{ch:introduction}

Micro-expressions are brief, involuntary facial movements that reveal an emotion which a person is attempting to conceal~\cite{yan2014}. They last no longer than approximately half a second~\cite{yan2014}, and their intensity is so low that an observer watching in real time will usually fail to notice them. A micro-expression may consist, for example, of a slight tightening of the brow or a movement at one corner of the mouth that is inconsistent with the expression the person is deliberately displaying.

This combination of informative content and near-invisibility has motivated sustained research into their automatic recognition. Yan et al.~\cite{yan2014} motivate the CASME II dataset on the grounds that a robust automatic recogniser ``would have broad applications in national safety, police interrogation, and clinical diagnosis''; the capability is likewise reported as promising for lie detection, business negotiation and psychoanalysis~\cite{xia2020a}. These are the motivating claims of the field rather than records of deployed systems, but they account for the continued research interest in the problem.

Automatic recognition systems for this task are typically assembled as pipelines of techniques that were introduced to the field separately. This thesis decomposes one such pipeline and measures the contribution of each of its stages individually.


\section{Research Motivation}
\label{sec:motivation}

Automatic recognition of micro-expressions from video cannot be achieved by applying a classifier directly to face images. Three properties of the phenomenon work against it. The facial movement is short, so only a small number of video frames carry the information that distinguishes one emotion from another. It is faint, so the information those frames carry lies close to the noise level of ordinary video. And it cannot be produced on instruction, so the recordings that contain it must be elicited under controlled laboratory conditions, which limits the resulting datasets to a small size.

Systems built for the task have responded to these constraints in a consistent way. The small facial motion is amplified, the sequence of pixel intensities is re-expressed as a description of movement, and the clip is resampled to a fixed duration before any learned component receives it; a spatial feature extractor then reduces each frame, and a temporal model processes the resulting sequence. Each stage was introduced with its own justification and its own reported gain, and each tends to be adopted together with the others.

This arrangement presents a practical difficulty for anyone building such a system. Implementation effort, and the tuning of preprocessing parameters that have no principled setting, are finite, so the developer must decide which stages justify that investment. The published literature does not answer this question, because each component is almost always evaluated within a complete system that differs from its predecessor in more than one respect.


\section{Research Aim and Objectives}
\label{sec:objectives}

The aim of this thesis is to investigate whether improved temporal modelling and motion representations can enhance the recognition of micro-expressions in dynamic facial expression recognition. The central research question is accordingly: \textit{how can the modelling of temporal phases and motion cues improve recognition accuracy for subtle and short-lived facial expressions?} The intended outcome is a clearer understanding of the limitations of current approaches, and of how sensitivity to fine-grained facial dynamics can be improved while robustness to noise is maintained.

The work builds on established supervised models and evaluates practical strategies for temporal and motion modelling. It does so by treating a standard micro-expression recognition pipeline as the object of study and measuring what each of its stages separately contributes under a single fixed protocol. Every configuration receives the same hybrid motion representation, in which each clip is encoded as horizontal optical flow, vertical optical flow and optical strain. Four components are then varied: Eulerian video magnification, which amplifies subtle facial motion before the motion representation is computed; a convolutional spatial stem, which learns spatial features from that representation; parameter-free attention (SimAM), which re-weights the learned features; and a temporal transformer, which models the sequence through temporal self-attention. The central research question is therefore addressed through three more specific questions:

\begin{enumerate}
	\item To what extent does temporal modelling with a transformer encoder improve recognition, compared with an otherwise identical pipeline without it?
	\item Does amplifying subtle facial motion through Eulerian video magnification, before the motion representation is computed, improve recognition?
	\item Do a learned convolutional spatial stem and parameter-free attention add to the recognition performance obtained from the hand-computed motion representation alone?
\end{enumerate}

To answer these questions, the thesis pursues the following objectives:

\begin{enumerate}
	\item \textbf{To implement a modular recognition pipeline} in which each of the four components can be enabled or disabled independently, with every disabled component replaced by a near-parameter-free fallback so that a comparison removes the component and nothing else.
	\item \textbf{To train and evaluate every architecturally valid configuration under one protocol.} The four components, switched on and off independently, give sixteen possible combinations, of which twelve are architecturally valid. Each of the twelve is trained from scratch, without pre-training or auxiliary data, which is the regime that a single small dataset permits. Each is evaluated under complete 25-fold leave-one-subject-out cross-validation on CASME II, using 156 clips in three classes of 99, 32 and 25 examples once the residual \texttt{Others} category is excluded. Because one fold is held out per subject, no subject's clips appear in both the training and the test data. Every factor not under variation is held identical across the twelve runs, and all twelve are scored by the same primary metric, pooled macro F1.
	\item \textbf{To estimate the marginal contribution of each component} from matched pairs of configurations. Any two configurations that differ in exactly one component isolate that component with all other factors fixed, so the difference between their scores estimates its contribution.
	\item \textbf{To examine the limitations of the resulting models}, by analysing per-class behaviour on the minority classes, comparing the results with published work, and stating the conditions under which the measured effects hold.
\end{enumerate}


\section{Research Contributions}
\label{sec:intro-contributions}

This thesis makes the following contributions.

\begin{enumerate}
	\item \textbf{A matched-pair factorial evaluation of a standard micro-expression recognition pipeline.} The reviewed literature typically evaluates components within complete systems, so that an improvement can rarely be attributed to the component credited with it. This thesis evaluates all twelve architecturally valid combinations of four components on CASME II under an identical leave-one-subject-out protocol, so that the contribution of each component, including learned ones, can be measured against an otherwise identical pipeline.

	\item \textbf{A measured account of the contribution of temporal and motion modelling.} Temporal modelling carries almost all of the measured effect: the temporal transformer improves pooled macro F1 by $+0.2173$ on average, and all six of its matched pairs are positive. The remaining components contribute far less. Eulerian video magnification has a small mean effect ($+0.0152$) whose direction varies across pairs; parameter-free attention has a mean effect ($+0.0034$) that cannot be distinguished from noise at this sample size; and the convolutional spatial stem is the only component with a negative mean effect ($-0.0310$).

	\item \textbf{Evidence that the full pipeline is not the best configuration.} A temporal transformer applied directly to the motion representation achieves a pooled macro F1 of $0.7122$, compared with $0.6659$ for the complete configuration in which all four components are enabled. An evaluation of the complete system alone would have reported the lower of these figures, with no means of discovering that a subset of its own architecture performs better.
\end{enumerate}


\section{Structure of the Thesis}
\label{sec:structure}

The thesis is organised in six chapters.

\begin{itemize}
	\item \textbf{Chapter 1: Introduction} motivates the study, states its aim, research questions and objectives, and summarises its contributions.

	\item \textbf{Chapter 2: Background} introduces the concepts on which the thesis builds: micro-expressions and the datasets used to study them, including CASME~II; data acquisition and preprocessing; Eulerian video magnification; optical flow and optical strain; neural network fundamentals and building blocks; training with limited and imbalanced data; and evaluation methodology.

	\item \textbf{Chapter 3: Literature Review} reviews the published evidence for each pipeline component and closes with the consolidated research gap and the declared divergences of the implemented system from the reviewed methods.

	\item \textbf{Chapter 4: Methodology} specifies the study as carried out: the experimental design, the dataset and its preparation, the twelve configurations, the architecture and its ablatable components, the training procedure and the evaluation protocol, together with the research gaps and limitations of each motion- and temporal-modelling component.

	\item \textbf{Chapter 5: Results} reports the results of the ablation, taking each component in turn over the matched pairs that isolate it, and then examines per-class behaviour, compares the results with published work and discusses them.

	\item \textbf{Chapter 6: Conclusion} answers the research questions and states the contributions, limitations and future work of the thesis.
\end{itemize}

%
% Cross-references used below (Structure of the Thesis only):
%   ch:literature            -> Chapter 3 (Literature Review)
%   sec:synthesis-and-gap    -> Section 3.10  (Synthesis and Gap)
%
% Labels defined here: ch:introduction, sec:motivation, sec:objectives,
%   sec:intro-contributions, sec:structure
%
% Bibliography keys: yan2014 (Yan et al., CASME II), xia2020a (Xia et al., STRCN)
%
% The aim, research question and objectives in Section 1.2 follow the third
% paragraph of the task description (contents/task_description.tex).
% The figures in Section 1.3 are taken from Chapter 5 (sec:transformer-effect,
% sec:evm-effect-result, sec:simam-effect-result, sec:cnn-effect-result,
% sec:headline-result) and must be kept consistent with it.
% =============================================================================

\chapter{Introduction}
\label{ch:introduction}

Micro-expressions are brief, involuntary facial movements that reveal an emotion which a person is attempting to conceal~\cite{yan2014}. They last no longer than approximately half a second~\cite{yan2014}, and their intensity is so low that an observer watching in real time will usually fail to notice them. A micro-expression may consist, for example, of a slight tightening of the brow or a movement at one corner of the mouth that is inconsistent with the expression the person is deliberately displaying.

This combination of informative content and near-invisibility has motivated sustained research into their automatic recognition. Yan et al.~\cite{yan2014} motivate the CASME II dataset on the grounds that a robust automatic recogniser ``would have broad applications in national safety, police interrogation, and clinical diagnosis''; the capability is likewise reported as promising for lie detection, business negotiation and psychoanalysis~\cite{xia2020a}. These are the motivating claims of the field rather than records of deployed systems, but they account for the continued research interest in the problem.

Automatic recognition systems for this task are typically assembled as pipelines of techniques that were introduced to the field separately. This thesis decomposes one such pipeline and measures the contribution of each of its stages individually.


\section{Research Motivation}
\label{sec:motivation}

Automatic recognition of micro-expressions from video cannot be achieved by applying a classifier directly to face images. Three properties of the phenomenon work against it. The facial movement is short, so only a small number of video frames carry the information that distinguishes one emotion from another. It is faint, so the information those frames carry lies close to the noise level of ordinary video. And it cannot be produced on instruction, so the recordings that contain it must be elicited under controlled laboratory conditions, which limits the resulting datasets to a small size.

Systems built for the task have responded to these constraints in a consistent way. The small facial motion is amplified, the sequence of pixel intensities is re-expressed as a description of movement, and the clip is resampled to a fixed duration before any learned component receives it; a spatial feature extractor then reduces each frame, and a temporal model processes the resulting sequence. Each stage was introduced with its own justification and its own reported gain, and each tends to be adopted together with the others.

This arrangement presents a practical difficulty for anyone building such a system. Implementation effort, and the tuning of preprocessing parameters that have no principled setting, are finite, so the developer must decide which stages justify that investment. The published literature does not answer this question, because each component is almost always evaluated within a complete system that differs from its predecessor in more than one respect.


\section{Research Aim and Objectives}
\label{sec:objectives}

The aim of this thesis is to investigate whether improved temporal modelling and motion representations can enhance the recognition of micro-expressions in dynamic facial expression recognition. The central research question is accordingly: \textit{how can the modelling of temporal phases and motion cues improve recognition accuracy for subtle and short-lived facial expressions?} The intended outcome is a clearer understanding of the limitations of current approaches, and of how sensitivity to fine-grained facial dynamics can be improved while robustness to noise is maintained.

The work builds on established supervised models and evaluates practical strategies for temporal and motion modelling. It does so by treating a standard micro-expression recognition pipeline as the object of study and measuring what each of its stages separately contributes under a single fixed protocol. Every configuration receives the same hybrid motion representation, in which each clip is encoded as horizontal optical flow, vertical optical flow and optical strain. Four components are then varied: Eulerian video magnification, which amplifies subtle facial motion before the motion representation is computed; a convolutional spatial stem, which learns spatial features from that representation; parameter-free attention (SimAM), which re-weights the learned features; and a temporal transformer, which models the sequence through temporal self-attention. The central research question is therefore addressed through three more specific questions:

\begin{enumerate}
	\item To what extent does temporal modelling with a transformer encoder improve recognition, compared with an otherwise identical pipeline without it?
	\item Does amplifying subtle facial motion through Eulerian video magnification, before the motion representation is computed, improve recognition?
	\item Do a learned convolutional spatial stem and parameter-free attention add to the recognition performance obtained from the hand-computed motion representation alone?
\end{enumerate}

To answer these questions, the thesis pursues the following objectives:

\begin{enumerate}
	\item \textbf{To implement a modular recognition pipeline} in which each of the four components can be enabled or disabled independently, with every disabled component replaced by a near-parameter-free fallback so that a comparison removes the component and nothing else.
	\item \textbf{To train and evaluate every architecturally valid configuration under one protocol.} The four components, switched on and off independently, give sixteen possible combinations, of which twelve are architecturally valid. Each of the twelve is trained from scratch, without pre-training or auxiliary data, which is the regime that a single small dataset permits. Each is evaluated under complete 25-fold leave-one-subject-out cross-validation on CASME II, using 156 clips in three classes of 99, 32 and 25 examples once the residual \texttt{Others} category is excluded. Because one fold is held out per subject, no subject's clips appear in both the training and the test data. Every factor not under variation is held identical across the twelve runs, and all twelve are scored by the same primary metric, pooled macro F1.
	\item \textbf{To estimate the marginal contribution of each component} from matched pairs of configurations. Any two configurations that differ in exactly one component isolate that component with all other factors fixed, so the difference between their scores estimates its contribution.
	\item \textbf{To examine the limitations of the resulting models}, by analysing per-class behaviour on the minority classes, comparing the results with published work, and stating the conditions under which the measured effects hold.
\end{enumerate}


\section{Research Contributions}
\label{sec:intro-contributions}

This thesis makes the following contributions.

\begin{enumerate}
	\item \textbf{A matched-pair factorial evaluation of a standard micro-expression recognition pipeline.} The reviewed literature typically evaluates components within complete systems, so that an improvement can rarely be attributed to the component credited with it. This thesis evaluates all twelve architecturally valid combinations of four components on CASME II under an identical leave-one-subject-out protocol, so that the contribution of each component, including learned ones, can be measured against an otherwise identical pipeline.

	\item \textbf{A measured account of the contribution of temporal and motion modelling.} Temporal modelling carries almost all of the measured effect: the temporal transformer improves pooled macro F1 by $+0.2173$ on average, and all six of its matched pairs are positive. The remaining components contribute far less. Eulerian video magnification has a small mean effect ($+0.0152$) whose direction varies across pairs; parameter-free attention has a mean effect ($+0.0034$) that cannot be distinguished from noise at this sample size; and the convolutional spatial stem is the only component with a negative mean effect ($-0.0310$).

	\item \textbf{Evidence that the full pipeline is not the best configuration.} A temporal transformer applied directly to the motion representation achieves a pooled macro F1 of $0.7122$, compared with $0.6659$ for the complete configuration in which all four components are enabled. An evaluation of the complete system alone would have reported the lower of these figures, with no means of discovering that a subset of its own architecture performs better.
\end{enumerate}


\section{Structure of the Thesis}
\label{sec:structure}

The thesis is organised in six chapters.

\begin{itemize}
	\item \textbf{Chapter 1: Introduction} motivates the study, states its aim, research questions and objectives, and summarises its contributions.

	\item \textbf{Chapter 2: Background} introduces the concepts on which the thesis builds: micro-expressions and the datasets used to study them, including CASME~II; data acquisition and preprocessing; Eulerian video magnification; optical flow and optical strain; neural network fundamentals and building blocks; training with limited and imbalanced data; and evaluation methodology.

	\item \textbf{Chapter 3: Literature Review} reviews the published evidence for each pipeline component and closes with the consolidated research gap and the declared divergences of the implemented system from the reviewed methods.

	\item \textbf{Chapter 4: Methodology} specifies the study as carried out: the experimental design, the dataset and its preparation, the twelve configurations, the architecture and its ablatable components, the training procedure and the evaluation protocol, together with the research gaps and limitations of each motion- and temporal-modelling component.

	\item \textbf{Chapter 5: Results} reports the results of the ablation, taking each component in turn over the matched pairs that isolate it, and then examines per-class behaviour, compares the results with published work and discusses them.

	\item \textbf{Chapter 6: Conclusion} answers the research questions and states the contributions, limitations and future work of the thesis.
\end{itemize}

%
% Cross-references used below (Structure of the Thesis only):
%   ch:literature            -> Chapter 3 (Literature Review)
%   sec:synthesis-and-gap    -> Section 3.10  (Synthesis and Gap)
%
% Labels defined here: ch:introduction, sec:motivation, sec:objectives,
%   sec:intro-contributions, sec:structure
%
% Bibliography keys: yan2014 (Yan et al., CASME II), xia2020a (Xia et al., STRCN)
%
% The aim, research question and objectives in Section 1.2 follow the third
% paragraph of the task description (contents/task_description.tex).
% The figures in Section 1.3 are taken from Chapter 5 (sec:transformer-effect,
% sec:evm-effect-result, sec:simam-effect-result, sec:cnn-effect-result,
% sec:headline-result) and must be kept consistent with it.
% =============================================================================

\chapter{Introduction}
\label{ch:introduction}

Micro-expressions are brief, involuntary facial movements that reveal an emotion which a person is attempting to conceal~\cite{yan2014}. They last no longer than approximately half a second~\cite{yan2014}, and their intensity is so low that an observer watching in real time will usually fail to notice them. A micro-expression may consist, for example, of a slight tightening of the brow or a movement at one corner of the mouth that is inconsistent with the expression the person is deliberately displaying.

This combination of informative content and near-invisibility has motivated sustained research into their automatic recognition. Yan et al.~\cite{yan2014} motivate the CASME II dataset on the grounds that a robust automatic recogniser ``would have broad applications in national safety, police interrogation, and clinical diagnosis''; the capability is likewise reported as promising for lie detection, business negotiation and psychoanalysis~\cite{xia2020a}. These are the motivating claims of the field rather than records of deployed systems, but they account for the continued research interest in the problem.

Automatic recognition systems for this task are typically assembled as pipelines of techniques that were introduced to the field separately. This thesis decomposes one such pipeline and measures the contribution of each of its stages individually.


\section{Research Motivation}
\label{sec:motivation}

Automatic recognition of micro-expressions from video cannot be achieved by applying a classifier directly to face images. Three properties of the phenomenon work against it. The facial movement is short, so only a small number of video frames carry the information that distinguishes one emotion from another. It is faint, so the information those frames carry lies close to the noise level of ordinary video. And it cannot be produced on instruction, so the recordings that contain it must be elicited under controlled laboratory conditions, which limits the resulting datasets to a small size.

Systems built for the task have responded to these constraints in a consistent way. The small facial motion is amplified, the sequence of pixel intensities is re-expressed as a description of movement, and the clip is resampled to a fixed duration before any learned component receives it; a spatial feature extractor then reduces each frame, and a temporal model processes the resulting sequence. Each stage was introduced with its own justification and its own reported gain, and each tends to be adopted together with the others.

This arrangement presents a practical difficulty for anyone building such a system. Implementation effort, and the tuning of preprocessing parameters that have no principled setting, are finite, so the developer must decide which stages justify that investment. The published literature does not answer this question, because each component is almost always evaluated within a complete system that differs from its predecessor in more than one respect.


\section{Research Aim and Objectives}
\label{sec:objectives}

The aim of this thesis is to investigate whether improved temporal modelling and motion representations can enhance the recognition of micro-expressions in dynamic facial expression recognition. The central research question is accordingly: \textit{how can the modelling of temporal phases and motion cues improve recognition accuracy for subtle and short-lived facial expressions?} The intended outcome is a clearer understanding of the limitations of current approaches, and of how sensitivity to fine-grained facial dynamics can be improved while robustness to noise is maintained.

The work builds on established supervised models and evaluates practical strategies for temporal and motion modelling. It does so by treating a standard micro-expression recognition pipeline as the object of study and measuring what each of its stages separately contributes under a single fixed protocol. Every configuration receives the same hybrid motion representation, in which each clip is encoded as horizontal optical flow, vertical optical flow and optical strain. Four components are then varied: Eulerian video magnification, which amplifies subtle facial motion before the motion representation is computed; a convolutional spatial stem, which learns spatial features from that representation; parameter-free attention (SimAM), which re-weights the learned features; and a temporal transformer, which models the sequence through temporal self-attention. The central research question is therefore addressed through three more specific questions:

\begin{enumerate}
	\item To what extent does temporal modelling with a transformer encoder improve recognition, compared with an otherwise identical pipeline without it?
	\item Does amplifying subtle facial motion through Eulerian video magnification, before the motion representation is computed, improve recognition?
	\item Do a learned convolutional spatial stem and parameter-free attention add to the recognition performance obtained from the hand-computed motion representation alone?
\end{enumerate}

To answer these questions, the thesis pursues the following objectives:

\begin{enumerate}
	\item \textbf{To implement a modular recognition pipeline} in which each of the four components can be enabled or disabled independently, with every disabled component replaced by a near-parameter-free fallback so that a comparison removes the component and nothing else.
	\item \textbf{To train and evaluate every architecturally valid configuration under one protocol.} The four components, switched on and off independently, give sixteen possible combinations, of which twelve are architecturally valid. Each of the twelve is trained from scratch, without pre-training or auxiliary data, which is the regime that a single small dataset permits. Each is evaluated under complete 25-fold leave-one-subject-out cross-validation on CASME II, using 156 clips in three classes of 99, 32 and 25 examples once the residual \texttt{Others} category is excluded. Because one fold is held out per subject, no subject's clips appear in both the training and the test data. Every factor not under variation is held identical across the twelve runs, and all twelve are scored by the same primary metric, pooled macro F1.
	\item \textbf{To estimate the marginal contribution of each component} from matched pairs of configurations. Any two configurations that differ in exactly one component isolate that component with all other factors fixed, so the difference between their scores estimates its contribution.
	\item \textbf{To examine the limitations of the resulting models}, by analysing per-class behaviour on the minority classes, comparing the results with published work, and stating the conditions under which the measured effects hold.
\end{enumerate}


\section{Research Contributions}
\label{sec:intro-contributions}

This thesis makes the following contributions.

\begin{enumerate}
	\item \textbf{A matched-pair factorial evaluation of a standard micro-expression recognition pipeline.} The reviewed literature typically evaluates components within complete systems, so that an improvement can rarely be attributed to the component credited with it. This thesis evaluates all twelve architecturally valid combinations of four components on CASME II under an identical leave-one-subject-out protocol, so that the contribution of each component, including learned ones, can be measured against an otherwise identical pipeline.

	\item \textbf{A measured account of the contribution of temporal and motion modelling.} Temporal modelling carries almost all of the measured effect: the temporal transformer improves pooled macro F1 by $+0.2173$ on average, and all six of its matched pairs are positive. The remaining components contribute far less. Eulerian video magnification has a small mean effect ($+0.0152$) whose direction varies across pairs; parameter-free attention has a mean effect ($+0.0034$) that cannot be distinguished from noise at this sample size; and the convolutional spatial stem is the only component with a negative mean effect ($-0.0310$).

	\item \textbf{Evidence that the full pipeline is not the best configuration.} A temporal transformer applied directly to the motion representation achieves a pooled macro F1 of $0.7122$, compared with $0.6659$ for the complete configuration in which all four components are enabled. An evaluation of the complete system alone would have reported the lower of these figures, with no means of discovering that a subset of its own architecture performs better.
\end{enumerate}


\section{Structure of the Thesis}
\label{sec:structure}

The thesis is organised in six chapters.

\begin{itemize}
	\item \textbf{Chapter 1: Introduction} motivates the study, states its aim, research questions and objectives, and summarises its contributions.

	\item \textbf{Chapter 2: Background} introduces the concepts on which the thesis builds: micro-expressions and the datasets used to study them, including CASME~II; data acquisition and preprocessing; Eulerian video magnification; optical flow and optical strain; neural network fundamentals and building blocks; training with limited and imbalanced data; and evaluation methodology.

	\item \textbf{Chapter 3: Literature Review} reviews the published evidence for each pipeline component and closes with the consolidated research gap and the declared divergences of the implemented system from the reviewed methods.

	\item \textbf{Chapter 4: Methodology} specifies the study as carried out: the experimental design, the dataset and its preparation, the twelve configurations, the architecture and its ablatable components, the training procedure and the evaluation protocol, together with the research gaps and limitations of each motion- and temporal-modelling component.

	\item \textbf{Chapter 5: Results} reports the results of the ablation, taking each component in turn over the matched pairs that isolate it, and then examines per-class behaviour, compares the results with published work and discusses them.

	\item \textbf{Chapter 6: Conclusion} answers the research questions and states the contributions, limitations and future work of the thesis.
\end{itemize}
